\section{Conclusion}\label{sec:conclusion}

In the fifteen months after the GENIUS Act became law, stablecoin inflows have not been associated with a larger fall in 3-month Treasury-bill yields than before. The point estimate is about half the pre-enactment one (\bPost{} against \bPre{}\,bp per USD~10\,bn at ten days), and the difference is imprecise (\bDiffP{}). The concentration of the association at the shortest maturities predates the Act. The outflow asymmetry seen before enactment is absent after it and is plausibly an artefact of the 2022--23 tightening cycle. A significant placebo break shows that the controls-only association is unstable, so even a significant change could not have been credited to the law.

None of this establishes that the Act has no effect on bill pricing. The rules are not yet binding, the post-enactment sample is short and the design is underpowered. Re-estimating the same pre-registered specification once the reserve rules have bound for a year, and with an instrument such as that of~\cite{ahmed2025}, would be the natural next step. The pipeline is built so that this requires only \texttt{make refresh}.
